% Generated by `taskgen course examples` from number_theory.inverse (seed 4).
% Do not edit: change examples.yaml and regenerate.
\begin{example}[необратимый вычет]\label{ex:euclid-noninvertible}
Найти элемент, обратный к $a = 169$ в кольце вычетов по модулю $n = 221$, или доказать, что его не существует.

\par\noindent\textit{Решение.} Обратный к $a$ существует тогда и только тогда, когда $\gcd(a, n) = 1$; расширенный алгоритм Евклида для $(n, a)$ даёт $d = \gcd(n, a)$ и $u n + v a = d$.\[\begin{array}{r|r|r|r}
r_i & q_i & u_i & v_i \\ \hline
221 &  & 1 & 0 \\
169 &  & 0 & 1 \\
52 & 1 & 1 & -1 \\
13 & 3 & -3 & 4 \\
0 & 4 & 13 & -17 \\
\end{array}\qquad r_i = u_i \cdot 221 + v_i \cdot 169\]
$d = \gcd(221, 169) = 13 > 1$. Если бы $a x \equiv 1 \pmod{n}$, то $ax - kn = 1$ для некоторого $k$, но левая часть делится на $13$ --- противоречие. Более того, $a$ --- делитель нуля: $169 \cdot 17 = 13 \cdot 221 \equiv 0\ (\textrm{mod}\ 221)$.\[\text{обратного элемента не существует}.\]
\end{example}
